\documentclass[11pt]{article}

\usepackage[margin=1in]{geometry}
\usepackage{amsmath}
\usepackage{amssymb}
\usepackage{booktabs}
\usepackage[numbers]{natbib}
\usepackage{hyperref}
\hypersetup{colorlinks=true, linkcolor=black, citecolor=black, urlcolor=blue}

\title{LatticeBench: A Contamination-Free Logic-Grid Benchmark for
Comparing Energy-Based Models, Language Models, and Constraint Solvers}

\author{David, Polished Snow}
\date{}

\begin{document}
\maketitle

\begin{abstract}
Static reasoning benchmarks decay as their instances leak into training data, and
the strongest formal-math suites are now largely solved. We introduce
LatticeBench, a procedurally-generated benchmark of logic-grid puzzles that is
distributed as a seed manifest rather than a fixed set of instances, so any
evaluator regenerates the exact puzzles bit-for-bit and no instance can be quietly
swapped or memorized without detection. Each puzzle has a unique solution verified
by a constraint solver and reduced to a minimal clue set, and each compiles to
three aligned representations: natural-language clues, a constraint program, and
an energy function over binary assignment variables whose only zero-energy
configuration is the solution. A single harness scores three model classes on the
identical instances: energy-based models, language models, and classical CP/SMT
solvers. This lets us report a graded energy metric that credits near-solutions,
place solvers as real baselines rather than difficulty oracles, and measure how
each class scales with grid size. All experimental numbers in this draft are
placeholders pending the full run.
\end{abstract}

\section{Introduction}

Benchmarks for machine reasoning wear out. A suite that separates models sharply
at release loses its power as its problems and solutions circulate, and the field
has watched formal-math benchmarks slide from single-digit pass rates toward
saturation within a few years \citep{minif2f}. The usual response is to build a
larger fixed set by hand, but hand-built sets inherit the same flaw on a slower
clock: they are finite, they cannot be regenerated, and they leak.

We take a different route. Instead of shipping puzzles, LatticeBench ships the
machine that makes them. A published manifest names a master seed, a generator
version, and a git commit, and from those three facts anyone reconstructs the same
batch of logic-grid puzzles bit-for-bit. Because the grader re-derives instances
from the manifest and re-grades server-side, a leaked or edited instance buys
nothing, and a fresh manifest is always one seed away.

Logic-grid puzzles are a good substrate for this because they are constraint
satisfaction problems with a single correct answer, they admit a difficulty dial
in the grid size, and they translate cleanly across representations. The same
puzzle can be read as clues by a language model, as a constraint program by a
solver, and as an energy function by an energy-based model. That alignment is what
lets us compare three model classes that are almost never measured on common
ground.

Our contributions are the following.

\begin{itemize}
  \item \textbf{Procedural, on-demand generation.} Instances are minted from a
  seed with grid size and clue budget as parameters, so difficulty is a
  continuous dial rather than a fixed distribution.
  \item \textbf{A rigorous contamination-free protocol.} We distribute a seed
  manifest, not solved puzzles, and enforce it with the steps below.
  \item \textbf{An energy formulation.} Every puzzle compiles to an energy
  function whose unique zero-energy configuration is the solution, making
  energy-based models a first-class evaluated class and yielding a graded
  final-energy metric with partial credit.
  \item \textbf{Solvers as real baselines.} Classical CP/SMT solvers are reported
  as a model class alongside language and energy models, not hidden as difficulty
  oracles.
\end{itemize}

The contamination protocol is enforced through four steps.

\begin{enumerate}
  \item Publish a manifest of master seed, generator version, and git commit.
  \item Regenerate instances deterministically and bit-for-bit from the manifest.
  \item Hash each instance back to its seed so substitution or editing is
  detectable.
  \item Date each manifest so held-out splits can be forward-dated past a model's
  training cutoff.
\end{enumerate}

\section{Related Work}

PutnamBench is the design we most closely emulate: it pairs multiple formal
representations with grading that factors the final answer from its justification,
reports heterogeneous baselines, and remains dramatically unsaturated behind a
public leaderboard \citep{putnambench}. Its limitation is that its problems are
hand-formalized and fixed, so the set cannot be regenerated and leaks over time.
MiniF2F illustrates where that leads, having moved from a hard release target
toward roughly 80 percent solved on its stronger splits \citep{minif2f}. The
closest procedural competitor is ZebraLogic, which grades logic-grid CSPs by
difficulty and documents an accuracy collapse as grids grow \citep{zebralogic};
it nonetheless ships a fixed generated pool, uses Z3 only as a difficulty oracle,
and evaluates neither solvers as baselines nor energy-based models. GridPuzzle is
a fixed set scraped from Puzzle Baron with no contamination control
\citep{gridpuzzle}, PuzzleBench pairs models with solvers as an execution backend
on a fixed set \citep{puzzlebench}, and the BIG-Bench logic-grid task is a small
static probe. Procedural natural-language reasoning suites such as ProofWriter
\citep{proofwriter} and PrOntoQA \citep{prontoqa} establish generation as accepted
methodology but target entailment rather than unique-assignment CSPs, while
PuzzleClone \citep{puzzleclone} and CombiBench \citep{combibench} extend puzzle and
combinatorics coverage without occupying the energy-and-solver corner. On the
energy side we build on the CSP-as-energy-landscape framing \citep{cspenergy},
self-supervised CSP solvers \citep{ssvcsp}, and Energy-Based Transformers
\citep{ebt}; Logical Intelligence's Kona EBRM and Aleph work is a non-peer-reviewed
blog source \citep{kona} and Enso is an open implementation \citep{enso}. Prior
work by the author on the Tars system reached 622/672 on PutnamBench, but that
figure comes from an LLM-orchestration pipeline (recursive decomposition, parallel
Lean, premise selection) in which an energy-based model acts only as a supporting
ranker; it is not a pure energy-based prover result and no LatticeBench number
derives from it.

Table~\ref{tab:diff} summarizes the positioning. LatticeBench is the only entry
that is procedural, contamination-free, evaluates EBMs as a class, reports solvers
as baselines, and provides a graded energy metric.

\begin{table}[t]
\centering
\begin{tabular}{lccccc}
\toprule
Benchmark & Procedural & Contam.-free & EBMs as class & Solvers baseline & Graded energy \\
\midrule
PutnamBench  & no  & no  & no  & no  & no  \\
MiniF2F      & no  & no  & no  & no  & no  \\
ZebraLogic   & yes & no  & no  & no  & no  \\
GridPuzzle   & no  & no  & no  & no  & no  \\
LatticeBench & yes & yes & yes & yes & yes \\
\bottomrule
\end{tabular}
\caption{Where LatticeBench sits relative to the closest benchmarks. Only
LatticeBench is ``yes'' in every column.}
\label{tab:diff}
\end{table}

\section{Benchmark Construction}

A LatticeBench instance places $n$ houses in a row and assigns $m$ attributes,
where each attribute has $n$ distinct values and each value occupies exactly one
house. The generator samples a ground-truth assignment, enumerates the clues from
a fixed vocabulary that are true of it, and adds clues one at a time until a
constraint solver confirms the solution is unique. It then greedily removes clues,
keeping a removal whenever the solver confirms uniqueness survives, so the final
clue set is minimal in the sense that no single clue can be dropped.

The clue vocabulary covers the standard logic-grid repertoire: equality,
negation, immediate and non-immediate positional relations, adjacency, ends of the
row, and ordinal position. Every clue type exposes both a solver constraint and an
energy penalty, which is what keeps the three representations aligned. The oracle
is CPMpy over OR-Tools CP-SAT with a Z3 fallback, and uniqueness is checked by
forbidding a found solution and re-solving.

Because all randomness flows through one seeded RNG and serialization is
canonical, the finished instance is a deterministic function of the seed and the
generator version. This determinism is what the contamination protocol depends on,
and it is verified in continuous integration by regenerating a reference batch and
comparing hashes across platforms.

\section{Contamination-Free Protocol}

We distribute a manifest, not puzzles. The manifest records a master seed, a
generator version, and the git commit that produced the release, together with the
batch configuration and a creation date. From those fields any evaluator
reconstructs the batch bit-for-bit, and the grader does the same server-side, so
it never trusts a submitted instance.

Each instance is defined as the generator's output on $(\text{master seed}, i,
\text{generator version})$ and carries a hash of its canonical form. A submitter
who swaps in an easier instance or edits a hard one changes the hash, and grading
recomputes the instances from the manifest regardless of what the submission
claims. Substitution and answer injection therefore fail by construction.

Creation dates handle the harder problem of memorization. A held-out split is a
manifest published on a known date; a model whose training cutoff precedes that
date cannot have seen instances that only exist by running the pinned generator on
the seed. We are explicit about the residual risk in Section~\ref{sec:limits}:
genre familiarity with zebra puzzles is real, legitimate, and distinct from
instance contamination, and we bound it by comparing a fresh forward-dated
manifest against an older public one of the same configuration.

\section{Energy Formulation}

We represent an assignment as binary variables $x_{a,v,h} \in \{0,1\}$, where
$x_{a,v,h}=1$ means attribute $a$ takes value $v$ in house $h$. The energy is

\begin{equation}
E(x) = \lambda \, S(x) + \sum_{c \in \mathcal{C}} P_c(x),
\end{equation}

where $S(x)$ is the structural term enforcing that each attribute takes exactly
one value per house and each value appears in exactly one house (one-hot and
all-different constraints), $\lambda > 0$ weights it, and each $P_c \geq 0$ is a
penalty for clue $c$ that is zero exactly when the clue is satisfied. Every term is
at most quadratic in $x$, so $E$ is a Potts/QUBO energy. Because the puzzle has a
unique solution and each penalty is nonnegative and vanishes only on satisfaction,
$E(x) = 0$ if and only if $x$ is that unique solution, and $E(x) > 0$ everywhere
else. An energy-based model is scored by minimizing $E$ and reading off the
argmin; the final value of $E$ gives a graded metric that credits configurations
close to the solution rather than only exact matches.

\section{Evaluation Setup}

The harness dispatches each instance to one of three adapters behind a common
interface. The solver adapter runs the oracle as a competitor and records its
conflict count and solve time. The language-model adapter renders clues to a
prompt and parses the returned grid. The energy-model adapter minimizes $E$ and
reads off the argmin configuration. Grading is factored in the PutnamBench style:
the final answer grid is scored separately from the justification, so a model that
reaches the right grid by faulty reasoning is not credited the same as one that
justifies it.

We report four metrics. Exact-match is the fraction of instances whose full grid
is correct. Cell-accuracy is the fraction of individual attribute-value-house
cells placed correctly, giving a smoother signal. Final-energy is the mean value
of $E$ at the model's chosen configuration, which is zero for a solved instance
and positive otherwise. Solve-time is wall-clock latency, reported chiefly for the
solver and energy classes where it is meaningful.

\section{Experiments}

This section describes the tables and plots the full run will produce; the numbers
in Table~\ref{tab:results} are placeholders. The headline plot is accuracy against
grid size for each model class, which we expect to show language-model accuracy
falling as grids grow while solvers stay near 100 percent, reproducing the
complexity curve ZebraLogic reported and setting it against a solver ceiling on the
same instances. A second analysis correlates the solver's conflict count with
language-model accuracy per instance, testing whether search difficulty for a
solver predicts difficulty for a language model. A third examines energy-versus-
correctness calibration, asking whether an energy-based model's final energy tracks
whether its answer is right.

\begin{table}[t]
\centering
\begin{tabular}{llcccc}
\toprule
Model & Class & Exact-match & Cell-acc & Mean energy & Median time (s) \\
\midrule
CP-SAT (OR-Tools) & Solver & TODO & TODO & TODO & TODO \\
Z3               & Solver & TODO & TODO & TODO & TODO \\
LLM baseline     & LLM    & TODO & TODO & TODO & TODO \\
EBM baseline     & EBM    & TODO & TODO & TODO & TODO \\
\bottomrule
\end{tabular}
\caption{Placeholder results. All numbers are TODO pending the full evaluation
run; the row structure fixes the reported model classes and metrics.}
\label{tab:results}
\end{table}

\section{Analysis}

The analysis will interpret the three studies above rather than restate them. We
are most interested in whether the difficulty orderings agree across classes: if
solver conflict count, language-model error, and energy-model residual energy rank
instances similarly, that argues logic-grid difficulty is a property of the
instance rather than of a particular solving method, which would make the grid-size
dial a clean difficulty control. Where the orderings disagree we expect the
disagreements to be informative about each method's failure mode.

\section{Limitations}
\label{sec:limits}

We state three limits plainly. First, energy-based models may underperform
language models on this benchmark; that is an informative result about the current
state of EBM reasoning, not a flaw in the benchmark, and reporting it honestly is
part of the point. Second, the benchmark is English-only, so it measures reasoning
entangled with English clue comprehension and does not speak to multilingual
performance. Third, models are broadly familiar with the zebra-puzzle genre, and
that familiarity raises scores legitimately; the protocol guarantees a specific
instance was unseen when a manifest postdates a model's cutoff, but it cannot make
a model naive about the genre, and we quantify the residual gap rather than deny
it.

\section{Conclusion}

LatticeBench replaces a fixed pool of reasoning problems with a generator, a seed
manifest, and a grader that re-derives everything server-side, which keeps the
benchmark useful after it is public and makes results hard to game. By compiling
each puzzle to aligned clue, constraint, and energy representations, it measures
energy-based models, language models, and constraint solvers on identical
instances and reports a graded energy metric alongside exact-match accuracy. The
experimental numbers here are placeholders; the contribution of this draft is the
construction, the protocol, and the evaluation design.

\begin{thebibliography}{99}

\bibitem{putnambench} G.~Tsoukalas et al. PutnamBench: Evaluating Neural
Theorem-Provers on the Putnam Mathematical Competition. arXiv:2407.11214, NeurIPS
Datasets and Benchmarks, 2024.

\bibitem{minif2f} K.~Zheng, J.~M.~Han, and S.~Polu. MiniF2F: A Cross-System
Benchmark for Formal Olympiad-Level Mathematics. arXiv:2109.00110, 2021.

\bibitem{zebralogic} B.~Y.~Lin et al. ZebraLogic: On the Scaling Limits of LLMs
for Logical Reasoning. arXiv:2502.01100, ICML, 2025.

\bibitem{gridpuzzle} Authors of GridPuzzle. GridPuzzle: A Benchmark for Evaluating
Reasoning Chains on Logic Grid Puzzles. arXiv:2407.14790, EMNLP, 2024.

\bibitem{puzzlebench} Authors of PuzzleBench. PuzzleBench: Can LLMs Solve Hard
Combinatorial Reasoning Puzzles? arXiv:2402.02611, 2024.

\bibitem{proofwriter} O.~Tafjord, B.~D.~Mishra, and P.~Clark. ProofWriter:
Generating Implications, Proofs, and Abductive Statements over Natural Language.
arXiv:2012.13048, 2020.

\bibitem{prontoqa} A.~Saparov and H.~He. Language Models Are Greedy Reasoners: A
Systematic Formal Analysis of Chain-of-Thought (PrOntoQA). arXiv:2210.01240, 2022.

\bibitem{puzzleclone} Authors of PuzzleClone. PuzzleClone: A Framework for Scaling
Formal Puzzle Benchmarks. arXiv:2508.15180, 2025.

\bibitem{combibench} Authors of CombiBench. CombiBench: Benchmarking LLM
Capability for Combinatorial Mathematics. arXiv:2505.03171, 2025.

\bibitem{cspenergy} Authors of the CSP-energy-landscape study. Constraint
Satisfaction as Energy-Landscape Search. arXiv:2501.00227, 2025.

\bibitem{ssvcsp} Authors of the self-supervised CSP solver. Self-Supervised
Learning for Constraint Satisfaction Problems. arXiv:2502.15794, 2025.

\bibitem{ebt} A.~Gladstone et al. Energy-Based Transformers Are Scalable Learners
and Thinkers. arXiv:2507.02092, 2025.

\bibitem{kona} Logical Intelligence. Kona: An Energy-Based Reasoning Model, and
the Aleph System. Company blog, 2025. Non-peer-reviewed.

\bibitem{enso} M.~V.~Pandey. Enso: Energy-Based Reasoning.
\url{https://github.com/MVPandey/Enso}, 2024.

\end{thebibliography}

\end{document}
